% Source: https://artofproblemsolving.com/community/c3h1982433_chocolate_math_problem
% Stable identity: aops-topic-1982433-q01
A chocolate bar is a $3\times4$ grid of unit squares. A \emph{break} takes one piece and splits it into two pieces along a single straight grid line that runs all the way across that piece. (So the first break produces a $1\times4$ and a $2\times4$ piece, or a $3\times1$ and a $3\times3$ piece, or a $3\times2$ and a $3\times2$ piece; a single unit square cannot be snapped off a corner.) Only one piece is broken at a time and pieces are not rotated or moved. The process continues until all pieces are unit squares.
\begin{enumerate}
\item[(a)] Show that exactly $11$ breaks are always used.
\item[(b)] Two processes are different if their sequences of breaks differ (a break is identified by the piece it splits and the line it uses). How many different processes are there?
\end{enumerate}
